\documentclass[10pt,a4paper]{amsart}
\usepackage[margin=19mm]{geometry}
\usepackage{amsmath,amssymb,mathtools}
\usepackage[scr=rsfs,cal=euler]{mathalfa}
\usepackage{xfrac}
\usepackage[unicode,hidelinks,bookmarks=false]{hyperref}
\newcommand{\ie}{i.e.\ }
\newcommand{\cf}{cf.\ }
\newcommand{\resp}{respectively\ }
\newcommand{\Subsection}{\subsection}
\providecommand{\nobreakdash}{\nobreak\hbox{-}}
\setlength{\emergencystretch}{2em}
\multlinegap0pt
\usepackage{amssymb}
\usepackage{xcolor}
\usepackage{tikz}
\newtheorem{theorem}{Theorem}
\title{Uniform stability of recovering the Sturm-Liouville operator on a star-graph}
\author{Maria Kuznetsova}
\date{Source: arXiv:2606.09405v1 (CC BY 4.0)}
\begin{document}
\maketitle

\noindent\emph{Source and licence.} Uniform stability of recovering the Sturm-Liouville operator on a star-graph. Version arXiv:2606.09405v1 (CC BY 4.0), distributed by arXiv under the Creative Commons Attribution 4.0 International licence, http://creativecommons.org/licenses/by/4.0/, as recorded in the arXiv licence metadata for this exact version and shown on the abstract page https://arxiv.org/abs/2606.09405v1. Full licence notice: \texttt{licenses/arxiv-2606.09405-CC-BY-4.0.txt}.\par\smallskip

\noindent\textit{Editorial scope.} Counted unit: the article's theorem stab (source0.tex lines 996--1008). The article's complete original proof of it is delivered with it (source0.tex lines 1184--1262, one \texttt{proof} environment of 79 lines, which the article places away from the statement). No mathematics is written, repaired or re-derived: the statement and the proof are the article's own lines, in the article's own notation, and no retained line is substituted or re-flowed.  Retained uncounted as the source-local context the proof needs: lines 1019--1025; lines 1027--1042; lines 1045--1047; lines 1074--1088; lines 1090--1092; lines 1166--1168.  Nothing was omitted from any retained span, so the package carries no omission record.  No retained line contains a citation, so the wrapper carries no bibliography and no bibliography entry.  No retained line contains a figure environment, an image-inclusion command or drawing code, so no image is delivered and the package declares no asset dependency.  Editorial formatting only, in addition: an AMS \texttt{amsart} preamble that restates the article's own declarations and macros used by the retained lines, namely the environments \texttt{theorem} on separate counters, together with the standard packages those lines need.  The retained environments print in this document as Theorem 1 in order of appearance, whereas the article prints the same items with the numbering of its own full document.  The complete native source of the article as distributed in the arXiv e-print for this exact version is bundled unchanged as \texttt{sources/202400/source0.tex}.
\begin{equation} \label{K_1}
P_1(x,t) = \frac{1}{2}Q\Big(\frac{x+t}{2}\Big) -
\frac{1}{2}Q\Big(\frac{x-t}{2}\Big),
\quad
K_1(x,t)=\frac{1}{2}Q\Big(\frac{x+t}{2}\Big)+\frac{1}{2}
Q\Big(\frac{x-t}{2}\Big),
\end{equation}

\begin{equation} \label{K_{n+1}}
\begin{split}
K_{n+1}(x,t)=\frac{1}{2}\int_t^x
\Bigg( \int_{\xi-t}^\xi q(\tau)K_n(\tau,t+\tau-\xi)\,d\tau +\int_{\frac{\xi+t}{2}}^\xi q(\tau)K_n(\tau,t-\tau+\xi)\,d\tau
\\
+\int_{\frac{\xi-t}{2}}^{\xi-t}
q(\tau)K_n(\tau,-t-\tau+\xi)\,d\tau \Bigg)\,d\xi,
\\
P_{n+1}(x,t)=\frac{1}{2}\int_t^x
\Bigg( \int_{\xi-t}^\xi q(\tau)P_n(\tau,t+\tau-\xi)\,d\tau +\int_{\frac{\xi+t}{2}}^\xi q(\tau)P_n(\tau,t-\tau+\xi)\,d\tau
\\
-
\int_{\frac{\xi-t}{2}}^{\xi-t}q(\tau)
P_n(\tau,-t-\tau+\xi)\,d\tau \Bigg)\,d\xi.
\end{split}
\end{equation}

\begin{equation}\label{K_n}
|P_n(x,t)|, \; |K_n(x,t)|\le \frac{x^{n-1}}{(n-1)!} Q^n_a(x), \quad n \ge 1, \quad Q_a(x) := \int_0^x |q(\xi)|\,d\xi.
\end{equation}

\begin{equation} \label{first der}
\begin{split}
K_1'(x,t) = \frac{1}{4} q\Big(\frac{x+t}{2}\Big) + \frac{1}{4}q\Big(\frac{x-t}{2}\Big),
\quad {\dot K}_1(x, t) = \frac14 q\Big(\frac{x+t}{2}\Big) - \frac14 q\Big(\frac{x-t}{2}\Big),
\\[2mm]
K_{n+1}'(x, t) = \frac12 \int_{x-t}^x q(\tau) K_n(\tau, t + \tau - x) \,d\tau + \frac12 \int_{\frac{x+t}{2}}^x q(\tau) K_n(\tau, t-\tau+x)\,d\tau \\
+\frac12 \int_{\frac{x-t}{2}}^{x-t} q(\tau) K_n(\tau,x-t-\tau)\,d\tau,
\\
{\dot K}_{n+1}(x, t) = -\frac{1}{2} \int_0^t q(\tau) K_n(\tau, \tau) \, d\tau +\frac12 \int_t^x \Bigg(\int_{\xi-t}^\xi q(\tau){\dot K}_n(\tau, t+\tau-\xi)\,d\tau \\
+\int_{\frac{\xi+t}{2}}^\xi q(\tau){\dot K}_n(\tau, t-\tau+\xi)\,d\tau 
- \int_{\frac{\xi-t}{2}}^{\xi-t} q(\tau){\dot K}_n(\tau, -t-\tau+\xi)\,d\tau\Bigg)\,d\xi \\
+ \frac{1}{4} \int_t^x q\Big(\frac{\xi-t}{2}\Big) K_n\Big(\frac{\xi-t}{2},\frac{\xi-t}{2}\Big) \,d\xi
- \frac{1}{4} \int_t^x q\Big(\frac{\xi+t}{2}\Big) K_n\Big(\frac{\xi+t}{2},\frac{\xi+t}{2}\Big)\, d\xi.
\end{split}
\end{equation}

\begin{equation} \label{K' est}
|K_{n}'(x, t)| \le \frac{x^{n-1}}{n!}Q_a^{n+1}(x), \quad |\dot K_{n}(x, t)| \le 2 \frac{x^{n-2}}{(n-1)!}Q^{n}_a(x), \quad 0 \le t \le x \le \pi.
\end{equation}

\begin{equation} \label{der series}
K'(x, t) = K'_1(x,t) + \sum_{n=2}^\infty K'_{n}(x, t), \quad \dot K(x, t) = \dot K_1(x,t) + \sum_{n=2}^\infty \dot K_{n}(x, t),
\end{equation}

\begin{theorem} \label{trans stab}
Let $R>0.$ Then, for any potentials $q$ and $\tilde q$ satisfying the condition
$$\| q\|_{L_2(0, \pi)} \le R, \quad \| \tilde q \|_{L_2(0, \pi)}\le R,$$
 the following estimates hold:
\begin{equation} \label{K stability}
\begin{split} 
\| \hat K'(x, \cdot)\|_{L_2(0, x)} \le B_R \| \hat q\|_{L_2(0, \pi)}, \quad \| \hat{\dot{K}}(x, \cdot)\|_{L_2(0, x)} \le B_R \| \hat q\|_{L_2(0, \pi)},
\\ 
\| \hat P'(x, \cdot)\|_{L_2(0, x)} \le B_R \| \hat q\|_{L_2(0, \pi)}, \quad \| \hat{\dot {P}}(x, \cdot)\|_{L_2(0, x)} \le B_R \| \hat q\|_{L_2(0, \pi)}, 
\end{split} 
\end{equation}
where $x \in (0, \pi]$ and $B_R = \frac{\sqrt 2}{2} + 12R \pi^{\frac32}\exp(2\pi^{\frac32}R).$
\end{theorem}

\begin{proof}[Proof of Theorem~\ref{trans stab}]
I. 
From formula~\eqref{K_1}, we obtain the inequality
\begin{equation*} \label{ind base}
|\hat K_1(x, t)| \le \hat Q_a(x) \le \sqrt{x} \| \hat q\|_{L_2(0, x)}, \quad \hat Q_a(x) := \int_0^x |\hat q(\tau)|\,d\tau.
\end{equation*}
From~\eqref{K_{n+1}}, it follows that
\begin{equation*}
\begin{split}
|\hat K_{n+1}(x, t)| \le \frac12 \int_t^x\Bigg( \int_{\xi-t}^\xi |\hat q(\tau)| |K_n(\tau,t+\tau-\xi)|\,d\tau +\int_{\frac{\xi+t}{2}}^\xi |\hat q(\tau)| |K_n(\tau,t-\tau+\xi)|\,d\tau \\
+\int_{\frac{\xi-t}{2}}^{\xi-t} |\hat q(\tau)| |K_n(\tau,-t-\tau+\xi)|\,d\tau +  \int_{\xi-t}^\xi |\tilde q(\tau)| |\hat K_n(\tau,t+\tau-\xi)|\,d\tau \\
+\int_{\frac{\xi+t}{2}}^\xi |\tilde q(\tau)||\hat K_n(\tau,t-\tau+\xi)|\,d\tau + \int_{\frac{\xi-t}{2}}^{\xi-t} |\tilde q(\tau)||\hat K_n(\tau,-t-\tau+\xi)|\,d\tau\Bigg) \,d\xi,
\end{split}
\end{equation*}
where $n \ge 1.$
Using these inequalities and~\eqref{K_n}, by induction, we get the estimates
\begin{equation} \label{hat K_n}
|\hat K_n(x, t)|\le 2 \hat Q_a(x) \frac{Q_{m}^{n-1}(x) x^{n-1}}{(n-1)!}, \; n \ge 1, \quad Q_{m}(x) := \int_0^x \max\{ |q(\xi)|, |\tilde q(\xi)|\}\, d\xi.
\end{equation}


II. Using~\eqref{K_n},~\eqref{first der}, and~\eqref{hat K_n}, we obtain
$$|\hat K'_{n+1}(x, t)| \le 3\hat Q_a(x) \frac{Q^{n}_{m}(x) x^{n-1}}{(n-1)!}, \quad n \ge 1.$$
From the formula for $K'_1$ in~\eqref{first der}, it follows that
$$\| \hat K'_1(x, \cdot) \|_{L_2(0,x)} \le \frac{\sqrt2}{2}\| \hat q\|_{L_2(0, x)}.$$
Using these inequalities and~\eqref{der series}, we estimate
\begin{equation*}\begin{split}
\|\hat K'(x, \cdot)\|_{L_2(0, x)} \le \big\| \hat K'_1(x, \cdot) \big\|_{L_2(0,x)} + \sum_{n=1}^\infty \big\| \hat K'_{n+1}(x, \cdot) \big\|_{L_2(0,x)} \\
\le \frac{\sqrt2}{2}\| \hat q\|_{L_2(0, x)}
+ 3\sqrt x \sum_{n=1}^\infty \hat Q_a(x) Q^n_{m}(x) \frac{x^{n-1}}{(n-1)!}\\
= \frac{\sqrt2}{2}\| \hat q\|_{L_2(0, x)} +
3\sqrt x \hat Q_a(x) Q_{m}(x) e^{x Q_m(x)}.
\end{split}
\end{equation*}
Taking into account the fact that
\begin{equation} \label{Q_a}
\hat Q_a(x) \le \sqrt{x} \| \hat q \|_{L_2(0, x)}, \quad Q_{m}(x) \le \sqrt x (\| q\|_{L_2(0, x)} + \| \tilde q\|_{L_2(0, x)}) \le 2 R\sqrt x,
\end{equation}
we arrive at the first inequality in~\eqref{K stability}.

III. From~\eqref{first der}, it follows that
\begin{equation} \label{final est}
\begin{split}
|\hat {\dot K}_{n+1}(x, t)|  \le  \int_0^{x} |\hat q(\tau)| |K_n(\tau, \tau)| \,d\tau 
+ \int_0^{x} |\tilde q(\tau)| |\hat K_n(\tau, \tau)| \,d\tau 
\\
+\frac12\int_t^x \Bigg(\int_{\xi-t}^\xi |\hat q(\tau)||\dot K_n(\tau, t+\tau-\xi)|\,d\tau +
\int_{\frac{\xi+t}{2}}^\xi |\hat q(\tau)||\dot K_n(\tau, t-\tau+\xi)|\,d\tau \\
+\int_{\frac{\xi-t}{2}}^{\xi-t} |\hat q(\tau)||\dot K_n(\tau, -t-\tau+\xi)|\,d\tau\Bigg)d\xi 
+\frac12\int_t^x \Bigg(\int_{\xi-t}^\xi |\tilde q(\tau)|\big|\hat{\dot K}_n(\tau, t+\tau-\xi)\big|\,d\tau \\
+\int_{\frac{\xi+t}{2}}^\xi |\tilde q(\tau)||\hat{\dot K}_n(\tau, t-\tau+\xi)|\,d\tau+\int_{\frac{\xi-t}{2}}^{\xi-t} |\tilde q(\tau)||\hat{\dot K}_n(\tau, -t-\tau+\xi)|\,d\tau\Bigg)d\xi.
\end{split}
\end{equation}
Let $n=1.$
Using the formula for $\dot K_1$ in~\eqref{first der}, applying~\eqref{K_n} and~\eqref{hat K_n}, we obtain
\begin{equation*} \label{base}
|\hat{\dot K}_2(x, t)| \le 6 \hat Q_a(x) Q_{m}(x).
\end{equation*}
We prove by induction the estimate
\begin{equation} \label{*}
|\hat{\dot K}_{n+1}(x, t)| \le 6 \hat{Q}_a(x) Q_{m}^{n}(x) \frac{x^{n-1}}{(n-1)!}, \quad n \ge 1.
\end{equation}
For $n=1,$ the formula is proved. For $n > 1,$ we estimate the right-hand side of~\eqref{final est} using~\eqref{K_n}, \eqref{K' est}, \eqref{hat K_n}, and~\eqref{*}. We obtain
$$ |\hat{\dot K}_{n+1}(x, t)| \le Q_{m}^n(x) \hat{Q}_a(x) \frac{x^{n-1}}{(n-1)!} \left(1 + \frac{10}{n}\right),$$
wherein for $n \ge 2,$ the value $1 + \frac{10}{n}$ does not exceed $6.$ Then, estimate~\eqref{*} is proved by induction. 

From the formula for $\dot{K}_1$ in~\eqref{first der},  it follows that
$$\| \hat {\dot K}_1(x, \cdot)\|_{L_2(0, x)} \le \frac{\sqrt 2}{2}\| \hat q\|_{L_2(0, x)}.$$
Applying~\eqref{der series},~\eqref{*}, and this inequality, we obtain
\begin{equation*}
\begin{split}
\|\hat{\dot K}(x, \cdot)\|_{L_2(0, x)} \le \big\| \hat{\dot K}_1(x, \cdot) \big\|_{L_2(0,x)} + \sum_{n=1}^\infty \big\| \hat{\dot K}_{n+1}(x, \cdot) \big\|_{L_2(0,x)}
\\
\le  \frac{\sqrt2}{2}\| \hat q\|_{L_2(0, x)} + 6\sqrt x \sum_{n=1}^\infty \hat Q_a(x) Q^n_{m}(x) \frac{x^{n-1}}{(n-1)!}\\
= \frac{\sqrt2}{2}\| \hat q\|_{L_2(0, x)} + 6\sqrt xQ_a(x) Q_{m}(x) e^{x Q_m(x)}.
\end{split}
\end{equation*}
Taking into account~\eqref{Q_a}, we arrive at the second inequality in~\eqref{K stability}.
\end{proof}

\end{document}
